\documentclass{article}
\usepackage[utf8]{inputenc}
\usepackage{amsmath}
\usepackage{amsfonts}
\usepackage{amssymb}
\usepackage{geometry}
\usepackage{tcolorbox}
\usepackage{hyperref}
\geometry{a4paper, margin=1in}

\title{Module 07: Optimization \& Advanced Ensembling - Part 2}
\author{Cybernauts 2026 Datathon Campaign}
\date{May 2026}

\begin{document}

\maketitle

\begin{abstract}
Ensembling represents the pinnacle of predictive performance in tabular data competitions. While individual models struggle with localized biases and structural blind spots, combining diverse architectures creates a collective intelligence that is robust to overfitting. This note explores the geometric and statistical mechanics of model blending and averaging, defining the conditions under which ensembling succeeds.
\end{abstract}

\tableofcontents
\newpage

\section{The Philosophy of Diverse Perspectives}
When competing at the highest levels of data science, relying on a single algorithm—no matter how meticulously tuned—creates a systemic vulnerability. Every machine learning algorithm interprets a feature space through a distinct mathematical lens.

\subsection{Algorithmic Blind Spots}
Consider the two dominant gradient boosting frameworks used in tabular hackathons:
\begin{itemize}
    \item \textbf{LightGBM:} Focuses on leaf-wise tree growth, optimizing for deep, asymmetric patterns. It handles high-order non-linear feature interactions exceptionally well but can occasionally overfit localized structural noise.
    \item \textbf{XGBoost:} Utilizes level-wise tree growth backed by strict L1 ($\alpha$) and L2 ($\lambda$) regularization penalties built directly into its formal objective function. This makes it highly stable against extreme outliers and distribution shifts.
\end{itemize}

An individual model will inevitably make structural mistakes on specific rows. However, because their optimization geometries differ, they rarely make the *same* mistakes on the *same* rows.

\section{The Mechanics of Blending and Averaging}
Model blending is the process of taking the raw prediction vectors or probability arrays from multiple distinct models and combining them using a defined mathematical operation.

\subsection{The Mathematics of Linear Combination}
For a classification task evaluated via log-loss or ROC-AUC, we extract the continuous prediction probabilities ($\hat{y}$) rather than hard-coded class predictions. A weighted linear combination ensemble is formalized as:
\begin{equation}
    \hat{y}_{\text{ensemble}} = \sum_{i=1}^{M} w_i \hat{y}_i \quad \text{subject to} \quad \sum_{i=1}^{M} w_i = 1
\end{equation}
Where $M$ is the number of individual models, and $w_i$ represents the operational weight assigned to model $i$.

\subsection{The Baseline Blend}
A standard, reliable starting configuration for a tabular ensembling pipeline pairs your optimized LightGBM and XGBoost models using a $60/40$ split:
$$\hat{y}_{\text{ensemble}} = 0.6 \cdot \hat{y}_{\text{LightGBM}} + 0.4 \cdot \hat{y}_{\text{XGBoost}}$$

This arithmetic formulation acts as a regularizing transformation. By smoothing the probability boundaries, it pulls erratic, overconfident predictions back toward the global mean, decreasing overall variance without inflating model bias.

\section{The Cardinal Rule: Uncorrelated Errors}
The success of an ensembling layer is strictly bound by a fundamental law of information theory:

\begin{tcolorbox}[colback=green!5 White,colframe=green!70!black,title=The Diversity Commandment]
Blending only yields performance improvements if the underlying constituent models make uncorrelated errors.
\end{tcolorbox}

\subsection{The Correlation Metric}
If you calculate the Pearson correlation coefficient ($r$) between the out-of-fold prediction vectors of Model A and Model B, you want to see a value significantly below $0.95$.
\begin{itemize}
    \item **High Correlation ($r > 0.98$):** Averaging two identical LightGBM models that use the same features and parameters provides zero predictive value. You are simply replicating the same error boundaries.
    \item **Low Correlation ($r < 0.90$):** Averaging a tree-based model with a deeply regularized linear model or a neural network almost always guarantees a performance bump up the leaderboard. The ensemble cancels out individual variances.
\end{itemize}

\newpage
\section{Practice Problems}

\textbf{P1:} You are building a binary classification ensemble. Model A outputs a prediction probability of $0.90$ for a true positive row, while Model B outputs a prediction probability of $0.40$ for the same row. Calculate the final ensemble prediction probability using a weighted blend of $0.7 \cdot \text{Model A} + 0.3 \cdot \text{Model B}$.
\\
\textit{Answer: Substitute the individual probabilities into the linear blending equation:
\begin{equation*}
    \hat{y}_{\text{ensemble}} = (0.7 \times 0.90) + (0.3 \times 0.40)
\end{equation*}
\begin{equation*}
    \hat{y}_{\text{ensemble}} = 0.63 + 0.12 = 0.75
\end{equation*}
The final ensemble probability is $0.75$.}

\newline
\textbf{P2:} A teammate wants to blend five distinct LightGBM models. Each model was trained on the exact same feature set, using the exact same cross-validation folds, varying only the random seed parameter. Explain why this approach violates the core philosophy of ensembling.
\\
\textit{Answer: This approach violates the principle of model diversity. Changing only the random seed results in minor structural variations during tree splitting, meaning the five models will remain highly correlated ($r \approx 0.99$). Because they share the same underlying bias and make nearly identical mistakes on the validation sets, averaging them fails to smooth out structural errors and adds unnecessary compute overhead to the pipeline.}

\newline
\textbf{P3:} In a regression datathon evaluated by Root Mean Squared Error (RMSE), you notice that your LightGBM model overpredicts values for high-income individuals, while a simple Ridge Regression model consistently underpredicts values for that same demographic. Explain why a linear blend of these two models will likely outperform either model individually.
\\
\textit{Answer: Because the models make opposite errors on the same demographic slice, their residuals are negatively correlated. Mathematically, when you compute the weighted average of the overpredicted value ($+\epsilon_1$) and the underpredicted value ($-\epsilon_2$), the opposing error vectors partially cancel each other out. This drags the ensemble's prediction closer to the true target line, significantly lowering the overall RMSE score.}

\end{document}
